% ======================================================================
% sections/introduction.tex
%
% Introduction scaffold for the paper:
%   Triple Humanoid 24/7 Adverse Event Oncology Trial Response Team:
%   4-Site Rotation
%
% This file is a draft scaffold. A future Claude Code Opus 4.7 1M Max
% session reads the bracketed instruction blocks below, then reads the
% files those instructions name, then expands the scaffold into the
% finished introduction. Target length is roughly 6 to 8 pages of the
% 70 plus page final paper.
% ======================================================================

\begin{sectionaccent}

[INTRODUCTION INSTRUCTIONS FOR FUTURE CLAUDE CODE OPUS 4.7 1M MAX SESSION.
Write the introduction as 5 connected subsections inside this single
\texttt{\textbackslash section\{Introduction\}} block. The 5 subsections are
1) Background And Setting, 2) Thesis Statement And Scope, 3) The
On-Premises Repository LLM Pattern, 4) The 4-Site Continental Rotation,
and 5) Roadmap Of The Paper. Use the bracketed file lists below to
draw the content. Add a verbatim ASCII diagram lifted from the file
named below into the 4-Site Continental Rotation subsection.

Files to read before writing Background And Setting:
\begin{itemize}
  \item \path{demo-projects/07-humanoid/paper/instructions/README.md}
        (whole file).
  \item \path{demo-projects/07-humanoid/paper/instructions/00-project-overview/}
        (whole directory, every file).
  \item \path{demo-projects/07-humanoid/paper/inputs/Deep-Research-A/chunk1_paper_text.md}
        and \path{chunk2_tables_references.md} for the Medical Full
        Automation A context.
  \item \path{demo-projects/07-humanoid/paper/inputs/Deep-Research-B/chunk_01_executive_summary_and_fully_automated_sponsor.md},
        \path{chunk_02_federated_learning_and_patient_prediction.md},
        \path{chunk_03_clinagent.md}, and
        \path{chunk_04_comparative_overview_and_implications.md}
        for the Medical Rapid Prototyping B context.
  \item Author's December 2025 paper
        \path{kawchak_2025_18029100} for the prior evidence that
        LLMs handle FDA adverse event data effectively. Cite the bib
        key \path{kawchak_2025_18029100} and quote the DOI
        \path{10.5281/zenodo.18029100}.
\end{itemize}

Files to read before writing Thesis Statement And Scope:
\begin{itemize}
  \item \path{demo-projects/07-humanoid/paper/instructions/README.md}
        Thesis section (verbatim quote of the thesis sentence is
        required in this subsection).
  \item \path{demo-projects/07-humanoid/paper/codegen/architecture.md}
        whole file.
\end{itemize}

Files to read before writing The On-Premises Repository LLM Pattern:
\begin{itemize}
  \item \path{demo-projects/07-humanoid/paper/codegen/src/llm_planner.py}
        for the 1 Hz planner loop.
  \item \path{demo-projects/07-humanoid/paper/codegen/src/broadcaster.py}
        for the fan out structure that delivers the same command
        payload to all 3 robots at a site at the same tick.
  \item \path{demo-projects/07-humanoid/paper/codegen/src/prompt_frozen.md}
        for the frozen prompt content.
  \item \path{demo-projects/07-humanoid/paper/codegen/src/comms_intellectual.py}
        and
        \path{demo-projects/07-humanoid/paper/codegen/src/comms_physical.py}
        for the dual channel (60 GHz UWB plus IR beacon plus shared
        compute fabric) that turns the per-site Claude Opus 4.7 1M
        instance into the on-premises repository LLM.
\end{itemize}

Files to read before writing The 4-Site Continental Rotation:
\begin{itemize}
  \item \path{demo-projects/07-humanoid/paper/codegen/diagrams/network_layout.txt}
        Include this ASCII diagram verbatim inside an \texttt{asciibox}
        environment as Figure 1. Caption it
        ``PAT-NET-001 continental layout with one Claude Opus 4.7 1M
        on-prem broadcaster per site and 3 Unitree H2 EDU robots per
        site.''
  \item \path{demo-projects/07-humanoid/paper/codegen/diagrams/swarm_dance.txt}
        Brief callout reference. Save the full inclusion for the
        Methods section.
  \item \path{demo-projects/07-humanoid/paper/codegen/src/cross_site_bus.py}
        for the de-identified cross site summary bus.
  \item \path{demo-projects/07-humanoid/paper/execution/diagrams/02_4site_execution_sequence.txt}
        for the 4 site execution sequence. Include this ASCII diagram
        verbatim inside an \texttt{asciibox} environment as Figure 2.
\end{itemize}

Files to read before writing Roadmap Of The Paper:
\begin{itemize}
  \item \path{demo-projects/07-humanoid/paper/draft-paper/main.tex}
        for the section ordering (Methods, Results, Discussion,
        Limitations and Future Work, Conclusions).
\end{itemize}

Required content cues for the introduction:
\begin{itemize}
  \item Open by framing the United States obligation to lead in patient
        safety, efficacy, and speed for oncology adverse event
        response. Stay respectful but advocate for accelerating the
        new field of Physical AI oncology trial surgical procedures.
  \item Frame the FDA RTCT 1 hour SLA as a real, accountable timer.
        Pair it with the camaraderie factor of 3 reduction in single
        robot error.
  \item Cite \path{kawchak_2025_18029100} once when introducing the
        FDA adverse event reporting evidence base.
  \item Present each milestone in the production chain as already
        checked off so the industry can build on top of it rather
        than rebuilding it from scratch.
  \item Use the ASCII diagrams listed above verbatim. Do not redraw.
\end{itemize}

Required real-life feasibility insight to include in this section:
``As of May 2026 the Unitree H2 EDU is shipping at the spec sheet level
named in the instructions README (39 DOF, 1.8 m, 70 kg, IP65 outdoor
rated). Pairing it with one on-prem Claude Opus 4.7 1M instance per
site is feasible on conventional high end x86 or Apple Silicon
servers. The 1 Hz broadcast cadence is well within current network and
inference latency budgets. What remains is the IRB, sponsor, and FDA
RTCT pilot enrollment, all of which the paper treats as part of the
forward roadmap rather than as solved.''

Do not exceed 8 pages. Use single dashes only. Use \S\ for the section
sign. Avoid stranding 1 or 2 word lines on their own.]

The introduction will set the FDA RTCT pilot context, name the 4-site
continental network, restate the thesis, and walk the reader through
the on-premises repository LLM pattern, the per-site swarm camaraderie,
and the roadmap that the rest of the paper follows.

\sectiontable{Introduction planning matrix}{%
Context & I1 & 1 & instr, inputs & FDA RTCT & 1 to 2 paragraphs & Cites kawchak & Ready \\
Thesis & I2 & 2 & instructions README & Quote thesis sentence & 1 paragraph & Verbatim quote & Ready \\
Pattern & I3 & 3 & llm planner & On-prem LLM & 2 paragraphs & 1 Hz cadence & Ready \\
Network & I4 & 4 & net layout & ASCII Figure 1 & 1 figure & 4 sites visible & Ready \\
Roadmap & I5 & 5 & main.tex & Section walk & 1 paragraph & Reader path & Ready \\
}

\end{sectionaccent}
